\section{Experimental Setup}
\label{sec:experiments}

\subsection{Datasets}
\label{sec:datasets}

\paragraph{Hallucination detection.} For long-form hallucination
detection (Table~\ref{tab:main} and Appendix
Table~\ref{tab:token_level})
we evaluate on three corpora that anchor this task in recent work:
RAGTruth \citep{niu2024ragtruth}, a retrieval-augmented
generation corpus with span-level hallucination annotations across
QA, summarization, and data-to-text prompts;
FactScore \citep{min2023}, atomic-fact-verified biography
generations; and LongFact \citep{wei2024}, long-form factuality
prompts paired with the SAFE atomic-fact verifier of the same paper.
For each corpus we follow the original protocol to obtain per-claim
support labels and broadcast them to per-token labels via the
overlap-resolution rule of
Appendix~\ref{ablation:overlap_rule}, which the token-level eval
uses to break ties when two claim spans cover the same token. The
per-claim labels also serve as the auxiliary supervision signal
$\tilde{y}_m$ introduced in Section~\ref{sec:notation}.

\paragraph{Uncertainty quantification.} For sequence-level
uncertainty quantification (Table~\ref{tab:polygraph}) we use the
LM-Polygraph benchmark
\citep{fadeeva2023lmpolygraph,vashurin2025lmpolygraph}, following
\citet{vashurin2025lmpolygraph} for subset selection, prompt
formatting, and few-shot example sourcing. Quality measures are
accuracy for short-form QA (MMLU, GSM8k), AlignScore for long-form
QA (TriviaQA, CoQA), and COMET for translation (WMT14 Fr-En, WMT19 De-En).

\paragraph{LLMs.} All experiments use three instruction-tuned
open-weight LLMs:
Mistral-7B-Instruct-v0.2 \citep{albertqjiang2023} (Mis-7B-Inst),
Llama-2-13B-chat-hf \citep{touvron2023} (LlaMa-13B-Inst), and
Llama-3.1-70B-Instruct \citep{touvron2023} (LlaMa-70B-Inst). For
each LLM we collect post-residual hidden states across all
transformer layers and all generated tokens, yielding the
hidden-state activations
$A \in \mathbb{R}^{L \times T \times D}$ that drive our method
(Section~\ref{sec:notation}).

\subsection{Baselines}
\label{sec:baselines}
We organize baselines into three row bands that match
Tables~\ref{tab:main} and~\ref{tab:polygraph}, with token-level
results in Appendix Table~\ref{tab:token_level}. All baselines are run through the
LM-Polygraph evaluation infrastructure
\citep{fadeeva2023lmpolygraph,vashurin2025lmpolygraph} on identical
generations and gold quality measures; sequence-level scores are
re-applied per sentence or per token when needed.

\paragraph{Confidence.} Score uncertainty from the model's
next-token distribution over generated tokens, or combine confidence
signals with attention-based features:
Maximum Sequence Probability (the LM-Polygraph default),
Perplexity \citep{fomicheva2020unsupervised},
Mean Token Entropy \citep{fomicheva2020unsupervised},
Self-Certainty \citep{kang2025selfcertainty},
TokenSAR \citep{duan2023},
Mean Pointwise Mutual Information \citep{takayama2019pmi},
BoostedProbSequence \citep{dinh2025boostedprob},
Contextualized Sequence Likelihood (CSL) \citep{lin2024csl},
Attention Score \citep{sriramanan2024attention}, and
Recurrent Attention-based Uncertainty Quantification (RAUQ)
\citep{vazhentsev2025rauq}.

\paragraph{Mechanistic.} Read model internals (hidden states or
activation maps) rather than the output distribution alone:
Fact-Probe \citep{han2025factprobe},
Feature-Gaps \citep{bakman2025featuregaps},
HaMI \citep{niu2025hami}, an MIL-based detector with
adaptive top-$k$ token selection and uncertainty-augmented
representations, and
Act-ViT \citep{barshalom2025}, the pooled-activation sequence
activation-map baseline.

\paragraph{Ours.} TPE is reported in two variants:
\emph{Encoder Only} -- the Stage-2 detector
(Section~\ref{sec:temporal_head}) used as a stand-alone per-window
scorer aggregated by top-$k$ MIL; and the full three-stage
\emph{TPE} pipeline (Section~\ref{sec:stage3}). Both are
trained with PRR as the validation criterion, following the
LM-Polygraph protocol.

\paragraph{Evaluation Measure.}
We report AUROC as the primary metric in
Table~\ref{tab:main}, Appendix Table~\ref{tab:token_level}, and
the Prediction Rejection Ratio
\citep{malinin2021uncertainty} for
Table~\ref{tab:polygraph}:
\begin{equation}
    \mathrm{PRR} = \frac{\mathrm{AUC}_{\mathrm{unc}} - \mathrm{AUC}_{\mathrm{rnd}}}{\mathrm{AUC}_{\mathrm{oracle}} - \mathrm{AUC}_{\mathrm{rnd}}},
    \label{eq:prr}
\end{equation}
which measures how well an uncertainty score ranks predictions by
quality relative to an oracle. We compute PRR up to a rejection
threshold of 50\%. AUROC, ECE, and Brier scores under the same
setup are reported in Appendix~\ref{app:extended}.

\paragraph{Generation setup.} All numbers are reported under
greedy decoding. Stochastic-sampling and best-sample variants
follow the LM-Polygraph defaults and are reported in
Appendix~\ref{app:extended}.
